\documentclass[12pt]{article}
\usepackage{latexsym,pptalk}

\font\colossal=cmss10 scaled 5000
\font\gigantic=cmss10 scaled 2500


%\transparency
%\onpaper

\def\pptalk{{
\raise.2ex\hbox{\textcolor{magenta}P}\kern-.25em
\raise-.1ex\hbox{\textcolor{red}p}\textcolor{blue}{talk}}} 
\def\ecomm#1{{\textcolor{blue}{$\backslash${\bf #1}}}}
\def\noncomm#1{{\textcolor{blue}{\bf #1}}}
\def\comm#1{{\textcolor{red}{$\backslash${\bf #1}}}}
\def\button#1{{\fboxrule=.5mm\framebox{#1}}}
\def\file#1{\textcolor{darkpurple}{\bf #1}}
\def\txtem#1{\texttt{\sl `#1'}}

\definecolor{beige}{rgb}{.9,.9,.7}
\definecolor{leaf}{rgb}{.75,1,.75}
\definecolor{darkpurple}{named}{RedViolet}


\title{\textcolor{blue}{\colossal The \pptalk\  package~: sample
and manual}} 
\author{\gigantic Palash B. Pal\\ 
\textcolor{red}{\gigantic pbpal@theory.saha.ernet.in}\\
Saha Institute of Nuclear Physics\\
1/AF Bidhan-Nagar, Calcutta 700064, INDIA\\[10mm]
}

\date{First release: February 2004\\ 
This version: March 2005}

\begin{document}



\maketitle

\beginslide

\section{Set it up first!}
%%
You are now viewing a PDF file.  I recommend that you see it with
\noncomm{Acroread} (called \noncomm{Acrobat Reader} in Windows), 5.0
or later.  \textcolor{red}{Please read the instructions on this slide
carefully before going to the next slide.}

%%
\begin{enumerate}

\item If you find that the file appears $90^\circ$ rotated as you open
it, use the \noncomm{View} menu to rotate it clockwise by
$90^\circ$. If the file appears in normal orientation, ignore this
step.

\item Use the \noncomm{full screen view} in the Acrobat.  This can be
accessed in the pull-down menu under the option \noncomm{View}, or
alternatively, by pressing \noncomm{Ctrl-L}.

\item Use the \button{left mouse} button, or the \button{PageDown}
button on the keyboard to go forward in the file.  In what follows,
this hit or click will be denoted by \button{$\bullet$}. 

\item The \button{right mouse} / \button{PageUp} buttons can be used
for going backward if you have to.

\item Even if you are not using \noncomm{full screen view} for some
reason, \underline{never use the scroll bar} to navigate within the
file.  Rather, click on the next-page arrow (which looks like
\button{$\rhd$}, a triangle pointing right) with your mouse.

\end{enumerate}


%%
Okay, now use \button{$\bullet$} to go to the next slide.

\endslide

\beginslide
\section{What \pptalk\ does}
%%
The package \pptalk\ is a \LaTeX-based style file which allows you to
achieve the following things:
%
\begin{enumerate}
\item Using a multi-media projector, you can project your talk on the
screen in different slides.  The slides can be divided into different
segments which can appear one at a time, at the push of one key on the
keyboard or one click of the mouse.


\item If you change your mind and decide to give the talk with
transparencies, you can use the same file, with one announcement in
the preamble, to generate a hard copy which will not divide a page
into segments.

\item If you further want to print out the talk where you would not
care about the division of the text into slides, you can again do it
with a different announcement in the preamble.

\end{enumerate}

\endslide

%%%%%%%%%%%%
\beginslide

\section{\pptalk\ vis-\`a-vis Power Point}
%
As you noticed, the objective is similar to the program {\em Power
Point\/} used in Windows OS.  Although I have never used {\em Power
Point\/}, I must admit that \pptalk\ has been inspired by {\em Power
Point\/} talk presentations that I have seen.  Indeed, some people
argue that the letters ``pp'' in the name \pptalk\ is merely an
acronym of this famous software, though some others say that they are
the initials of the person who wrote \pptalk.

Like {\em Power Point\/}, segments appear with clicks in \pptalk.
However, it does not have animation like {\em Power Point\/}.  On the
other hand, you can use the rich formatting features of \LaTeX,
including equations and tables etc.  If you are used to using \LaTeX,
you can use all its features in \pptalk.

\endslide


%%%%%%%%%%%%
\beginslide

\section{\pptalk\ vis-\`a-vis other \LaTeX-based presentation packages}
%
I have seen two more \LaTeX-based talk presentation packages:
\file{prosper} and \file{pdfscreen}.  I have not used either, and have
come to know about the latter one only after writing most of \pptalk.
Both these packages are much larger than \pptalk, and they can do many
things that \pptalk\ cannot.  However, there is a flip side to it.
Both these packages have many commands as opposed to only a few for
\pptalk.  \pptalk\ consists of only a style file, a few kilobytes
long, which can be easily transported and installed on another system
if you are away on a trip.  The other packages have to use separate
class files, which means that the input file looks very different from
what you would write for an ordinary \noncomm{article} class file.  In
short, \pptalk, though rudimentary, has the most important features
needed for a talk presentation, and is easier to deal with.  This, by
the way, is not the only advantage of \pptalk.  The user is left to
judge for others.

\endslide

%%%%%%%%%%%%%

\slidebox{red}{beige}

\beginslide

\section{How to use \pptalk}
%
Create the \LaTeX\ file following the instructions given later in this
document.  Then use the standard commands to produce the dvi file, the
ps file and the pdf file.  For example, if the input file is called
\file{xyz.tex}, issue the following commands in successsion:
%
\begin{quote}
\noncomm{latex\quad xyz\\
dvips\quad xyz\quad  -o\\
ps2pdf\quad  xyz.ps}
\end{quote}
%
The final stage produces a file called \file{xyz.pdf}.  You can also
create this file by using the alternative commands
%
\begin{quote}
\noncomm{latex\quad xyz\\
dvipdf\quad xyz.dvi\quad  $>$ \quad xyz.pdf}
\end{quote}
%
Open the pdf file on the screen using any standard pdf reader like
acroread (or Acrobat Reader).  As emphasized earlier, Acrobat 5.0 or
later versions should be used.


If you find that the file appears $90^\circ$ rotated as you open
it, use the VIEW menu to rotate it clockwise by $90^\circ$. If the
file appears in normal orientation, ignore this suggestion.

Now use the full screen view in the Acrobat, as discussed earlier.

\endslide

%%%%%%%%%%%%%

\beginslide

\section{How to organize the file}
%
You don't need to use a new class file in order to write your talk
with \pptalk.   You can use the \noncomm{article} class file.  Use the
package \pptalk.  In other words, start your file as follows:
%
\begin{quote}
\ecomm{documentclass[12pt]\{article\}} \\ 
\ecomm{usepackage\{pptalk\}}
\end{quote}
%
and of course, with the commands for beginning and ending the
document.  In the middle, you can write \LaTeX\ in the usual way, with
a few special commands for \pptalk\ which are discussed in the rest of
this document.  These special commands will appear in
\textcolor{red}{red} in this document.  Any other command, either a
shell command or commands belonging to standard \LaTeX\ or to any
existing style file, will appear in \textcolor{blue}{blue}.

\textcolor{blue}{Be careful to include the \file{12pt} option in the
\ecomm{documentclass} declaration, as shown above.  The sizes of fonts
have been set with this option in such a way that appears nice on the
projected screen.}

\underline{Do not use} any of the page style commands like
\ecomm{textheight} or \ecomm{textwidth}.  These have already been set
by \pptalk\ to some optimum values which will suit a screen
presentation.

\endslide

%%%%%%%%%%%%%%

\beginslide

\section{Dividing things into slides}
%
You need to decide how much (or how little) of text you want to put in
a particular slide.  When you want to begin a particular slide, type
%
\begin{quote}
\comm{beginslide}
\end{quote}
%
Then continue with whatever you want to put there, typing things just
as in a normal \LaTeX\ file.  

Divide things into paragraphs as you want, put equations as you want:
%
\begin{eqnarray*}
\vec x = \pmatrix{x_1 \cr x_2} \,.
%\label{}
\end{eqnarray*}
%
When you finish typing the contents of the slide, type
%
\begin{quote}
\comm{endslide}
\end{quote}
%
After that,  use the \comm{beginslide} command to start a new slide.

If you have ended up putting too little material on the slide, you
will obtain a warning message while running \LaTeX.  You may choose to
ignore this message, or want to put in some more material on the slide.

\endslide


%%%%%%%%%%%%%%
\beginslide
%
On the other hand, if the material is a bit too much to fit on a
slide, \pptalk\ will mildly complain.  While normally the page header
(containing the slide number) appears above the slide box, \pptalk\
will now put the header on an otherwise empty page, followed by a page
containing the slide box.  This means that you need to roughly shorten
the material by one line in order to make it fit.

When the material becomes so long that even that does not work,
\pptalk\ will give you a page with an error message.  Now some more
drastic changes are necessary.  Here are your options:
%%
\begin{enumerate}
\item You can transfer some of the contents of the slide to another
slide.  This is the best option.

\item You can reduce the font size.  The default font size is
\ecomm{LARGE}.  You can change it to, e.g., \ecomm{Large} by issuing the
command
%
\begin{quote}
\ecomm{def{\comm{defaultsize\{$\backslash$Large\}}}}
\end{quote}
%
anywhere within the file.  This is not recommended unless you are
giving the talk to a small audience, in a small room.

\item In principle, you can increase the size of the slides by putting
suitable values of \ecomm{textheight} and/or \ecomm{textwidth} in the
preamble, like you can do in any \LaTeX\ file.  This is highly
discouraged since the resulting slide may not fit on a computer
screen.

\end{enumerate}



\endslide



%%%%%%%%%%%%%%
\beginslide
\section{Dividing things into segments}
%
Now come the novel things.  Suppose you want to present the contents
of a page in a number of segments, one segment appearing at a time.
What do you do?

You decide how the segmentation will have to be made in a slide.  
Start writing the slide with the first segment, after the
\comm{beginslide} command.  When one segment is
done and you want to start the next segment, just type the command
%
\begin{quote}
\comm{newsegment}
\end{quote}
%
Then write the text for the next segment.  Continue this process until
you reach the end of the slide with \comm{endslide}.

Let us now see an example.  Notice that there is an empty portion at
the bottom of this slide.  This is because there are more segments
down there.  Use \button{$\bullet$} now.

\newsegment

\textcolor{magenta}{Now you see the next segment.}  You can continue this
way, introducing more segments.  In principle, the number of segments
on a page can be unlimited.  However, the length of the page and the
size of the font introduces some obvious restrictions.

From now on, hit \button{$\bullet$} whenever you want to see some more
material on the screen: whether material continued on the same slide
or on a different slide.

\endslide

%%%%%%%%%%%%%%%%

\beginslide

You can put the command \comm{newsegment} almost anywhere in the file.
You can put it between two paragraphs, as you have done once earlier
and will do again at the end of this paragraph.  You will realize that
the command can be put at the end of a sentence if you hit
\button{$\bullet$} now. \newsegment You can put the command before or
after an equation, before or after a table, or even in the middle of a
word like (let's see if I can spell it right) \textcolor{red}{ \Huge
``super\-\newsegment calli\newsegment fragi\newsegment
listic\newsegment expi\newsegment alli\newsegment docious''.}

%\newsegment
\endslide

%%%%%%%%%%%%%%%

\beginslide

You can also put the command within groups.  By \noncomm{group}, I
mean anything which is between a matching beginning and ending
command.  For example, it can be a table, an environment like quote or
enumerate, an equation, a command in the math mode which uses matching
brackets.

However, in this case you will have to do a bit of extra work.  The
command \comm{newsegment} works only within a group.  So, when the
group ends, if you want to continue with the same segment, put the
command 
%
\begin{quote}
\comm{samesegment}
\end{quote}
%
and keep writing.  For example, in writing
%
\begin{eqnarray*}
x =  \newsegment 
\textcolor{blue}{\pmatrix{x_1 \cr x_2} \,,}
%\label{}
\end{eqnarray*}
%
\samesegment 
I had to put the \comm{samesegment} command right after
the equation.  And when I wrote  
%
\begin{eqnarray*}
x = \left[ \frac1a + \newsegment \frac1b \right] \samesegment (y+2) 
%\label{}
\end{eqnarray*}
%
\samesegment using \ecomm{left[}$\cdots$\ecomm{right]}, 
I had to put the \comm{samesegment} command after
the closing square bracket.

\endslide

%%%%%%%%%%%%%%

\slidebox{red}{leaf}

\beginslide
\section{How to use colors}
%
You have already seen colors being used in this sample file.  This has
been done by using the package \file{color}.  \pptalk\ itself has this
package as an input, so you should not introduce it explicitly when
you are using \pptalk.

\newsegment

The \file{color} package defines two commands for implementing colored
text.  One of them is \ecomm{textcolor}, used for a local change of
color.  The syntax is
%
\begin{quote}
\ecomm{textcolor\{{\em colorname}\}\{Text material\}}
\end{quote}
%
which makes the \txtem{`Text material'} to appear in the color given
as \txtem{`colorname'}.  This command has been suitably redefined in
\pptalk, and will work exactly the same way within \pptalk.  On the
other hand, you can use the command
%
\begin{quote}
\comm{Color\{\em colorname\}}
\end{quote}
%
which will change the color of the text for the rest of the document.
This is a \pptalk\ special, and is a variation of the command
\ecomm{color} defined in the \file{color} package.


\endslide

%%%%%%%%%%%%%

\beginslide

\section{Color names}
%%
The basic color names appear in the file \file{color.sty}, which is
part of the color package.  They are: red, blue, green, cyan, magenta,
yellow, and the non-colors white and black.

\newsegment

Additional colors, defined by the file \file{color.pro} in the color
package, can also be used.  Take, e.g., the name ``Plum'' from this
file.  If you issue the command
%
\begin{quote}
\ecomm{definecolor\{{\em plum}\}\{named\}\{Plum\}} 
\end{quote}
%
you will be able to use \txtem{plum} for a \txtem{colorname}.

\newsegment

You can also create new colors.  For example, a new color called {\bf
RONG} can be defined by any of the following commands, where each
$p_i$ is a number between 0 and 1:
%
\begin{quote}%\label{p:defcol}
\ecomm{definecolor\{{\em RONG}\}\{rgb\}\{$p_1,p_2,p_3$\}} \\
\ecomm{definecolor\{{\em RONG}\}\{cmyk\}\{$p_1,p_2,p_3,p_4$\}} \\
\ecomm{definecolor\{{\em RONG}\}\{gray\}\{$p_1$\}}
\end{quote}
%
The last command produces only shades of gray.  Test with various
values of $p_i$ to create your own pallette.


\endslide

%%%%%%%%%%%%%%% 
\beginslide 

\section{Background color} 
%%
There can be two different meanings of the words ``background color''
in the context of \pptalk.  Any slide appears as a box on a page.  If
you want to change the color of the page outside the box, use the
standard command defined in the color package:
\begin{quote}
\ecomm{pagecolor\{{\em colorname}\}}
\end{quote}
%
where \txtem{colorname} is the name of a color already defined,
either in the file \file{color.sty} or by yourself.  The default color
is white, which is what has been used in this document.  You can
change it anywhere within the document.  The command takes effect from
the slide on which it appears.

You might also want to change the color of the box containing text,
and the color of the border.  For this, you need to use the \pptalk\
command 
%
\begin{quote}
\comm{slidebox\{{\em color1}\}\{{\em color2}\}}
\end{quote}
%
before any \comm{beginslide} command.  
Once you use this, \txtem{color1} will be the color of the
border, and \txtem{color2} will 
be the base color within the borders, starting from the slide which
starts after this declaration.  If you don't like the border, 
\txtem{color1} must be the same as the background color of the
page, or the base color set by \txtem{color2}.


\endslide

%%%%%%%%%%

\slidebox[1mm]{blue}{beige}
\beginslide

If you also want to change the width of the border, you can use an
optional argument.  For example, before this page, the border has been
set to blue, the base color to beige, and the borderwidth to 1mm by
issuing the command
%
\begin{quote}
\comm{slidebox[1mm]\{blue\}\{beige\}}
\end{quote}
%
where the color beige has been defined by
%
\begin{quote}
\ecomm{definecolor\{beige\}\{rgb\}\{.9,.9,.7\}}
\end{quote}
%
earlier in the document.  Note that the base color was
beige earlier as well.  But it does not matter.  You have to supply
the names of both colors in the \comm{slidebox} command even if you
want to change only one of them, or even if you want to just change
the borderwidth without changing the colors at all.

\newsegment

This command can also be utilized if you do not want the border around
the text.  Just set the width to zero.   Remember that the width must
have a unit of length.   For example, you can write
%
\begin{quote}
\comm{slidebox[0mm]\{blue\}\{beige\}}
\end{quote}
%
Instead of \txtem{mm}, you can use any other unit recognized by \LaTeX.



\endslide


%%%%%%%%%%

\beginslide

%%
I have a few recommendations on the use of colors.
%%
\begin{enumerate}
%
\item Do not use the global color change command \ecomm{color} defined
in the package \file{color.sty}.  Use \ecomm{textcolor} or
\comm{Color} instead.

\item Use your judgement for color contrasts.  Needless to say, the
contents of a slide will not be very legible if you use orange letters
on a yellow background.

\item Use light and dull shades for background.  Typeset most of the
things in black.  No other color appears as beautifully on the screen.
Use other colors only for emphasis.

\item It is best not to define a new color unless you absolutely need
to do so.  Use, as much as possible, the names defined in
\file{color.pro} in order to retain some kind of universality in the
color names.  It contains a lot of names, as shown next.  Try them
out.

\end{enumerate}
%%

\endslide
%%%%%%%%%%%%%%%
\beginslide

\begin{center}
\bf
\begin{tabular}{llll}
\hline
\multicolumn{4}{c}{Colors defined in \file{color.pro}} \\
\hline
GreenYellow & Yellow & Goldenrod & Dandelion \\ 
Apricot & Peach & Melon & YellowOrange \\ 
Orange & BurntOrange & Bittersweet & RedOrange \\ 
Mahogany & Maroon & BrickRed & Red \\ 
OrangeRed & RubineRed & WildStrawberry & Salmon \\ 
CarnationPink & Magenta & VioletRed & Rhodamine \\ 
Mulberry & RedViolet & Fuchsia & Lavender \\ 
Thistle & Orchid & DarkOrchid & Purple \\ 
Plum & Violet & RoyalPurple & BlueViolet \\ 
Periwinkle & CadetBlue & CornflowerBlue & MidnightBlue \\ 
NavyBlue & RoyalBlue & Blue & Cerulean \\ 
Cyan & ProcessBlue & SkyBlue & Turquoise \\ 
TealBlue & Aquamarine & BlueGreen & Emerald \\ 
JungleGreen & SeaGreen & Green & ForestGreen \\ 
PineGreen & LimeGreen & YellowGreen & SpringGreen \\ 
OliveGreen & RawSienna & Sepia & Brown \\ 
Tan & Gray & Black & White \\
\hline
\end{tabular} 
\end{center}
\endslide  

%%%%%%%%%%

\beginslide

\section{The running head on a slide}

Each slide contains a header.  For example, the header of the present
slide contains a line which says \txtem{Created with pptalk}, followed
by the slide number.  This is the default header, and will appear on
all slides unless you change it. 

To change the text appearing in the header to, say, \txtem{My favorite
header}, you need to issue the
command 
%
\begin{quote}
\ecomm{def\comm{runningheadtext\{\em My favorite
header\}}}
\end{quote}
%
In addition, you can also change the color of the running header to
\txtem{color1} by issuing the command
%
\begin{quote}
\ecomm{def\comm{runningheadcolor\{\em color1\}}}
\end{quote}
%
You can use these commands at the beginning of the file, or between
any two slides.

The slide number will appear automatically.  You need not do anything
for it.  Its color will be the same as the color of the text in the
header. 

\endslide

%%%%%%%%%%%%%%%
\slidebox[1mm]{blue}{leaf}

\beginslide
\section{Some technical points}
%%
\begin{itemize}

\item On a segmented slide, do not use objects which are automatically
numbered by \LaTeX.  The segmentation will interfere with the
numbering.  If you have to start a section, you can do it with the
usual \ecomm{section} command, which I have redefined to appear
without a number.  If you want to have a displayed equation, use the
\noncomm{eqnarray*} or the \noncomm{equation*} environments (with the
asterisk marks), or begin and end equations with the double dollar
(\$\$) signs, which will prevent the numbering of equations.

\item The contents of a slide are put into a vbox by \pptalk.
Environments like \noncomm{enumerate}, \noncomm{quote}, and
\noncomm{itemize} cannot work across different vboxes, and therefore
cannot be continued on different slides.  For long lists, it is
advisable not to use \noncomm{enumerate} at all.  Instead, use
\noncomm{itemize}.  That way, you can put a number of items on one
slide, end the \noncomm{itemize} environment on that slide, and
continue with a new \noncomm{itemize} on the next slide.  Looking at
the result, no one will know that the same \noncomm{itemize} is not
continuing.

\end{itemize}

\endslide

%%%%%%%%%%%%%%

\beginslide
%%
\section{Displaying tables}
\begin{itemize}
%%
\item Do not use the \noncomm{table} environment that \LaTeX\
provides.  Apart from putting a number of the table, it also puts the
table as a floating object, and \LaTeX\ decides where to put it.  Just
for creating the table, the \noncomm{tabular} environment is good enough.
If the entries contain too many math symbols, you can also use the
\noncomm{array} environment, putting the table within an equation.

\item Do not use the \comm{newsegment} command in the middle of a 
table.  I haven't been able to make it work.  This is a bug, which
will be fixed shortly.  Meanwhile, if you really want to highlight
something in a table, use colors to do that.

\end{itemize}

\endslide

%%%%%%%%%%%%%%%

\beginslide
%%
\section{Displaying figures}
\begin{itemize}
\item There are a number of packages for inserting postscript or EPSF
files into a \LaTeX\ file.  In \pptalk, I use \file{epsf.sty}, which
must not be introduced by the user because it is already included in
\pptalk.  

\item  For reasons described in the context of the \noncomm{table}
environment,  the \noncomm{figure} environment should not be used.  
Just insert your figure by the usual command \ecomm{epsfbox}, with one
specification mentioned below.

\item The normal syntax for \ecomm{epsfbox} is as follows:
%%
\begin{quote}
\noncomm{\ecomm{epsfxsize={\em len1}} \ecomm{epsfysize={\em len2}}
\ecomm{epsfbox\{{\em figfile}\}}}
\end{quote}
%%
where \noncomm{\em len1, len2} are lengths in any units recognized by
\LaTeX, and \noncomm{\em figfile} is a ps or eps file for a figure.
The size declarations are normally optional.  However, for \pptalk,
you must have an \ecomm{epsfysize} declaration.  In fact, \pptalk\
decides on the size of the slide from this declaration.  The
\ecomm{epsfxsize} declaration still remains optional.

\end{itemize}

\endslide

%%%%%%%%%%%%%%%

\beginslide
%%
\begin{itemize}
\item If you prefer some other figure package, you can still use it,
but then you will have to announce it by the \ecomm{usepackage}
command explicitly.  And I guess that the figures will then be
incompatible with the segmentation commands of \pptalk.  This means
that you will have to put the figures in a page which is not divided
into segments.

\item But if you use the \file{epsf} package, you can use segmentation
commands on pages containing figures.  An example is given on this
page.


\item On a segmented slide, avoid using very large figure files which
will take a long time to load.

\end{itemize}

\newsegment

\centerline{\epsfysize=8cm\epsfbox{pic.ps}}

\endslide

%%%%%%%%%%%%%%%

\beginslide
\section{On defaults}
%%
The default font family for \LaTeX\ is \noncomm{cmr}.  These fonts
do not appear well in pdf conversions.  So I have reset the default to
\noncomm{times}.  The \noncomm{palatino} and the \noncomm{helvetica}
families also work quite well.  If you want to use \noncomm{palatino},
for example, put the command
%
\begin{quote}
\ecomm{usepackage\{palatino\}}
\end{quote}
%
in the preamble of the file after the \pptalk\ package has been
included.  For \noncomm{helvetica}, the corresponding package is
called \noncomm{helvetic}.

Other defaults are summarized in the following table:
%%
\begin{center}
\begin{tabular}{lll}
\hline
Parameter & Default value & Can be changed by \\ 
\hline
Font size & \ecomm{LARGE} & \comm{defaultsize} \\
Text color & black & \comm{Color} \\
Background color & white & \ecomm{pagecolor} \\ 
Slide base color & cyan & \comm{slidebox} \\
Slide border color & red & \comm{slidebox} \\
Slide border width & 2mm & \comm{slidebox} \\
\hline
\end{tabular}
\end{center}


\endslide




%%%%%%%%%%

\beginslide

\section{Presentation tips not related to \pptalk}

\begin{list}{\textcolor{red}{$\spadesuit$}}{}

\item Acrobat reader has some special modes for appearance of a slide
  or a segment which appear in the full-screen view mode.  For
  example, in the PREFERENCES menu, you can set up such that new
  slides appear in the DISSOLVE mode.  Test to see what it does.

\item The \LaTeX\ package called \file{pstricks} has many beautiful
  features which can aid your talk presentation.  For example, you can
  draw an arrow from one place in the slide to another.

\item You can use a variety of symbols to announce an \comm{item} in a
  list like this one.  For example, the list on this  slide has been
  initiated by writing
%
\begin{quote}
\ecomm{begin\{list\}\{\ecomm{textcolor}\{red\}%
\{\$\ecomm{spadesuit}\$\}\}\{\}} 
\end{quote}
%
followed by the usual \ecomm{item} command and the command to end the
list.  Some fancy symbols for item markers are available in different
packages.

\end{list}


\endslide

%%%%%%%%%%%%%%%

\beginslide
\section{Printing out}
%
In case you want to print out the talk, you would presumably want all
segments of a page to appear simultaneously.  Two options have been
provided for this.  If you type
%
\begin{quote}
\comm{transparency}
\end{quote}
%
in the preamble of the file, you will obtain a version which you can
use on transparencies.  On the other hand, if you type 
%
\begin{quote}
\comm{onpaper}
\end{quote}
%
in the preamble, you will obtain a more condensed version, where text
will not be divided into slides, but will rather appear continuously,
as is usual for printing out an article.  You cannot use both at the
same time.  If you put both commands, only the last one will take
effect. 

If later you want to produce a multimedia talk again, just remove all 
of these declarations from the preamble.

\endslide

%%%%%%%%%%%%%%%

\beginslide
\section{Summary of commands}
%
Here is a summary of the special commands for \pptalk.
%
\begin{list}{}{}
\item[\comm{beginslide}$\,\cdots\,$\comm{endslide}] Denotes the
beginning and end of a slide.  The contents of the slide goes in the
middle, in place of the three dots.

\item[\comm{newsegment}] Denotes the beginning of a new segment in a
slide.  Not required for the first segment in a slide, which begins
with the \comm{beginslide} declaration.

\item[\comm{samesegment}] If the \comm{newsegment} appears within a
group, this command has to be used at the end of the group to reassert
the continuation of the segment after the group.

\item[\comm{slidebox}] Defines the border and base colors of all
slides which start after this announcement, unless changed by the same
command again.  An optional argument defines the width of the border.

\item[\comm{Color}] Changes the color of the text for the rest of the
document.

\item[\comm{transparency}] To be put in the preamble if you 
want all segments to appear simultaneously on a page, as e.g.\ is
usual for a transparency talk.

\item[\comm{onpaper}] To be put in the preamble if you want a
continuous printout, like you would want for an article to be sent to
a journal.

\end{list}
%%

\endslide

\beginslide

Then there are commands which need to be set with a leading
\ecomm{def}.  These are:
%
\begin{list}{}{}
\item[\ecomm{def}\comm{defaultsize}] Redefines the default size of
  letters in the text.

\item[\ecomm{def}\comm{runningheadtext}] Redefines the running head on
  the slides.

\item[\ecomm{def}\comm{runningheadcolor}] Redefines the color of the
  letters in the running head on the slides.

\end{list}
%%

Also useful are the following commands,  which
are not \pptalk\ specials but are not a standard \LaTeX\ commands 
either.  The following two are parts of the package \file{color}:
%
\begin{list}{}{}
\item[\ecomm{definecolor}] Used for defining colors.

\item[\ecomm{textcolor}] Redefines the text color locally.

\end{list}
%%

And then the command from the \file{epsf} pacakge:
%
\begin{list}{}{}
\item[\ecomm{epsfbox}] Used for inserting figures. Remember that the
\ecomm{epsfysize} command is mandatory here.

\end{list}
%%

\endslide



%%%%%%%%%%

\beginslide
%
\section{Acknowledgements}  I thank Soumitro Banerjee and Amitabha
Lahiri, who made many important comments and refused to be satisfied
by earlier versions of \pptalk.  I hope they are still dissatisfied,
which will provide me with motivations for further improvements of
\pptalk.

If you find any bug, or have any suggestion for improvements without
substantially complicating the program, please contact me.

\endslide


\end{document}

